% Computer Science A --- TypeScript (practice bank).
% Formatted from raw notes; Angular items are in \texttt{web-angular.tex}.

\runningheader{Computer Science A}{Web: TypeScript}{Page \thepage\ of \numpages}

\begingradingrange{csawebadvts}

\uplevel{\par\textbf{TypeScript}\par}

\question[4]
TRUE or FALSE: Generics allow reusable components that work with multiple types while maintaining type safety.
\begin{choices}
  \CorrectChoice TRUE
  \choice FALSE
\end{choices}
\begin{solution}
  Generics parameterize types so one function, class, or interface can serve many concrete types without falling back to \lstinline[style=tsinline]|any|. The checker verifies that arguments and return values stay consistent for each type argument.
\begin{lstlisting}[style=ts]
function identity<T>(value: T): T {
  return value;
}
const n = identity(42);       // T inferred as number
const s = identity('hello');  // T inferred as string
\end{lstlisting}
\end{solution}

\question[4]
TRUE or FALSE: TypeScript is a superset of JavaScript, meaning all valid JavaScript code is also valid TypeScript code.
\begin{choices}
  \CorrectChoice TRUE
  \choice FALSE
\end{choices}
\begin{solution}
  TypeScript builds on JavaScript syntax and runtime behavior. Existing \texttt{.js} files can often be renamed to \texttt{.ts} and adopted gradually as types are added. A few edge cases exist (for example very old or non-standard syntax), but the superset relationship is the guiding model for migration.
\begin{lstlisting}[style=js]
// Valid JavaScript is valid TypeScript (types optional)
function greet(name) {
  return 'Hello, ' + name;
}
\end{lstlisting}
\end{solution}

\question[4]
TRUE or FALSE: A tuple is a dynamic array that can hold any number of elements of different types.
\begin{choices}
  \choice TRUE
  \CorrectChoice FALSE
\end{choices}
\begin{solution}
  A \textbf{tuple} fixes both \emph{length} and \emph{position} of each element's type. It is not a growable heterogeneous array. \lstinline[style=tsinline]|string[]| or \lstinline[style=tsinline]|Array<string | number>| model variable-length collections.
\begin{lstlisting}[style=ts]
let user: [string, number] = ['Ada', 42];
// user.push(true);  // compile error: wrong length/type
// user = ['Ada'];   // compile error: missing second element
\end{lstlisting}
\end{solution}

\question[4]
TRUE or FALSE: A union type allows a variable to hold one of several types, written with the \lstinline[style=tsinline]!|! symbol.
\begin{choices}
  \CorrectChoice TRUE
  \choice FALSE
\end{choices}
\begin{solution}
  Union types combine alternatives with the pipe (\lstinline[style=tsinline]!|!) operator. The value must be assignable to at least one member of the union; narrowing (\lstinline[style=tsinline]|typeof|, \lstinline[style=tsinline]|instanceof|, type guards) is required before calling type-specific methods.
\begin{lstlisting}[style=ts]
let id: string | number;
id = 'a1';
id = 99;
// id.toFixed();  // error until narrowed with typeof id === 'number'
\end{lstlisting}
\end{solution}

\question[4]
Which compiler option is most important for preventing common \lstinline[style=tsinline]|null| and \lstinline[style=tsinline]|undefined| bugs? \emph{Select all that apply.}
\begin{checkboxes}
  \CorrectChoice \texttt{strict}
  \choice \texttt{target}
  \choice \texttt{module}
  \choice \texttt{outDir}
\end{checkboxes}
\begin{solution}
  \texttt{strict} enables a family of checks, including \textbf{strictNullChecks}: \lstinline[style=tsinline]|null| and \lstinline[style=tsinline]|undefined| are distinct types and cannot be assigned to \lstinline[style=tsinline]|string|, \lstinline[style=tsinline]|number|, and so on unless the type explicitly allows them (\lstinline[style=tsinline]|string | null|). \texttt{target} selects emitted ECMAScript version; \texttt{module} selects the module format; \texttt{outDir} only sets the output folder.
\begin{lstlisting}[style=ts]
// tsconfig.json excerpt
{
  "compilerOptions": {
    "strict": true
  }
}
\end{lstlisting}
\end{solution}

\question[4]
What is the return type of a function that loops forever and never returns normally? \emph{Select all that apply.}
\begin{checkboxes}
  \choice \lstinline[style=tsinline]|any|
  \choice \lstinline[style=tsinline]|undefined|
  \choice \lstinline[style=tsinline]|void|
  \CorrectChoice \lstinline[style=tsinline]|never|
\end{checkboxes}
\begin{solution}
  \lstinline[style=tsinline]|never| means no value is ever produced---the function does not complete (\lstinline[style=tsinline]|while (true) {}|) or always throws. \lstinline[style=tsinline]|void| applies to functions that \emph{do} finish without a meaningful return value. \lstinline[style=tsinline]|undefined| is a value, not the type of an unreachable endpoint.
\begin{lstlisting}[style=ts]
function spin(): never {
  while (true) { /* no return */ }
}
\end{lstlisting}
\end{solution}

\question[4]
What is a main difference between a \lstinline[style=tsinline]|type| alias and an \lstinline[style=tsinline]|interface|? \emph{Select all that apply.}
\begin{checkboxes}
  \CorrectChoice Interfaces support declaration merging; type aliases do not
  \choice Type aliases work only in browsers; interfaces work only in Node.js
  \choice Interfaces can describe only primitives; types describe only objects
  \choice Type aliases cannot represent union types, but interfaces can
\end{checkboxes}
\begin{solution}
  Declaring two interfaces with the same name in one scope \textbf{merges} their members. A duplicate \lstinline[style=tsinline]|type| name is a compile error. Both constructs are erased at compile time (no runtime presence in either browsers or Node.js). \textbf{Type aliases} can express unions and tuples; interfaces describe object shapes and support \lstinline[style=tsinline]|extends|.
\begin{lstlisting}[style=ts]
interface User { name: string; }
interface User { age: number; }
// Merged: User requires name and age

type Status = 'idle' | 'loading' | 'done';  // union via type alias
\end{lstlisting}
\end{solution}

\question[4]
TRUE or FALSE: If a variable is declared with type inference (\lstinline[style=tsinline]|let x = 10|), a string can later be assigned without a compile error.
\begin{choices}
  \choice TRUE
  \CorrectChoice FALSE
\end{choices}
\begin{solution}
  Inference locks the binding to the deduced type (\lstinline[style=tsinline]|number| here). Reassigning \lstinline[style=tsinline]|x = 'hello'| fails with \emph{Type 'string' is not assignable to type 'number'}. To allow both, annotate a union explicitly.
\begin{lstlisting}[style=ts]
let x = 10;
x = 'hello'; // compile error

let y: string | number = 10;
y = 'hello'; // valid
\end{lstlisting}
\end{solution}

\question[4]
TRUE or FALSE: The \lstinline[style=tsinline]|any| type is safer than \lstinline[style=tsinline]|unknown| because \lstinline[style=tsinline]|any| forces type narrowing before use.
\begin{choices}
  \choice TRUE
  \CorrectChoice FALSE
\end{choices}
\begin{solution}
  \lstinline[style=tsinline]|unknown| is the type-safe top type: most operations require narrowing first. \lstinline[style=tsinline]|any| disables checking and allows arbitrary property access and method calls, which can hide bugs until runtime.
\begin{lstlisting}[style=ts]
let value: unknown = JSON.parse('{"ok":true}');
// value.ok;  // error: object is unknown

if (typeof value === 'object' && value && 'ok' in value) {
  // narrowed safely
}
\end{lstlisting}
\end{solution}

\question[4]
TRUE or FALSE: Interfaces support declaration merging, so multiple declarations with the same name combine their properties.
\begin{choices}
  \CorrectChoice TRUE
  \choice FALSE
\end{choices}
\begin{solution}
  Merging is useful for extending third-party or global types without editing library source. Each declaration contributes properties to one merged interface.
\begin{lstlisting}[style=ts]
interface Window {
  myAppVersion: string;
}
// Augments the global Window interface when included in the project

interface User { name: string; }
interface User { age: number; }

const dev: User = { name: 'Jaehoon', age: 29 };
\end{lstlisting}
\end{solution}

\question[4]
TRUE or FALSE: An optional function parameter (marked with \lstinline[style=tsinline]|?|) must appear after all required parameters.
\begin{choices}
  \CorrectChoice TRUE
  \choice FALSE
\end{choices}
\begin{solution}
  Optional parameters may be omitted at call sites; placing them before required parameters would make argument positions ambiguous. Required parameters must come first.
\begin{lstlisting}[style=ts]
function greet(firstName: string, lastName?: string) {
  return lastName ? `${firstName} ${lastName}` : firstName;
}

// function bad(lastName?: string, firstName: string) {} // compile error
\end{lstlisting}
\end{solution}

\question[4]
TRUE or FALSE: In TypeScript, the \lstinline[style=tsinline]|number| type represents both integers and floating-point values.
\begin{choices}
  \CorrectChoice TRUE
  \choice FALSE
\end{choices}
\begin{solution}
  TypeScript follows JavaScript's IEEE-754 model: there is no separate \lstinline[style=tsinline]|int| or \lstinline[style=tsinline]|float| keyword. \lstinline[style=tsinline]|number| covers whole numbers, fractions, \lstinline[style=tsinline]|NaN|, and \lstinline[style=tsinline]|Infinity|. Very large integers may use \lstinline[style=tsinline]|bigint| (\lstinline[style=tsinline]|100n|).
\begin{lstlisting}[style=ts]
let count: number = 42;
let ratio: number = 3.14;
let invalid: number = NaN;
\end{lstlisting}
\end{solution}

\question[4]
TRUE or FALSE: Type assertion (\lstinline[style=tsinline]|as| or angle-bracket syntax) converts data from one type to another at runtime (for example, string to number).
\begin{choices}
  \choice TRUE
  \CorrectChoice FALSE
\end{choices}
\begin{solution}
  Assertions are compile-time hints only; they emit no conversion code. Asserting \lstinline[style=tsinline]|'123' as number| does not change the runtime string. Use \lstinline[style=tsinline]|Number()|, \lstinline[style=tsinline]|parseInt|, or similar for real conversion.
\begin{lstlisting}[style=ts]
let raw = '123';

let unsafe = raw as unknown as number; // still a string at runtime
let safe = Number(raw);              // actually 123 (number)
\end{lstlisting}
\end{solution}

\endgradingrange{csawebadvts}
